% ID: FIS-URT-QDM-001
% Titolo: Lanci successivi su una piattaforma priva di attrito
% Disciplina: Fisica
% Area: Quantita di moto e urti
% Argomento: Quantita di moto
% Sottoargomento: Rinculo e confronto tra sequenze di lancio
% Classe: Terza liceo scientifico
% Difficolta: 4
% Tipo: problema_numerico
% Risultato: T_ABC = 45,5 s; T_CBA = 45,2 s; velocita finale = 0,143 m/s in entrambi i casi
% Tempo_stimato: 15 min
% Competenze: conservazione della quantita di moto, moto rettilineo uniforme a tratti, confronto tra strategie
% Prerequisiti: quantita di moto, sistemi isolati, moto rettilineo uniforme
% Tag: fisica, quantita_moto, rinculo, sistema_isolato, moto_a_tratti
% Fonte: rielaborazione_da_traccia_fornita_dal_docente
% Autore: Riccardo Carli
% Licenza: CC-BY-SA-4.0

\begin{esercizio}
Un uomo di massa \(M=70\,\mathrm{kg}\) si trova inizialmente fermo al centro
di una piattaforma circolare di ghiaccio di raggio \(r=5{,}00\,\mathrm{m}\),
sulla quale puo scivolare senza attrito. Porta con se tre palle di masse
\[
m_A=1{,}00\,\mathrm{kg},\qquad
m_B=2{,}00\,\mathrm{kg},\qquad
m_C=3{,}00\,\mathrm{kg}.
\]
Le velocita indicate di seguito sono misurate rispetto alla piattaforma.

Nel primo scenario l'uomo lancia, sempre nella stessa direzione, la palla
\(A\) al tempo \(t=0\) con velocita \(3{,}00\,\mathrm{m/s}\), la palla \(B\)
dopo \(10{,}0\,\mathrm{s}\) con velocita \(2{,}00\,\mathrm{m/s}\) e la palla
\(C\) dopo altri \(10{,}0\,\mathrm{s}\) con velocita
\(1{,}00\,\mathrm{m/s}\).

Nel secondo scenario i lanci avvengono agli stessi istanti e con le stesse
velocita, ma nell'ordine inverso \(C\)-\(B\)-\(A\).

Determina:
\begin{enumerate}
\item in quale scenario l'uomo raggiunge prima il bordo della piattaforma;
\item se la sua velocita finale al bordo dipende dall'ordine dei lanci.
\end{enumerate}
\end{esercizio}

\begin{soluzione}
Scegliamo positivo il verso dei lanci. L'uomo si muove quindi nel verso
opposto. Poiche non agiscono forze esterne orizzontali, la quantita di moto
totale rimane nulla.

\textbf{Scenario \(A\)-\(B\)-\(C\).}

Dopo il primo lancio l'uomo porta ancora con se \(B\) e \(C\), quindi
\[
m_Av_A+(M+m_B+m_C)u_1=0.
\]
Da cui
\[
u_1=-\frac{1\cdot3}{70+2+3}
=-0{,}0400\,\mathrm{m/s}.
\]
Nei primi \(10{,}0\,\mathrm{s}\) percorre in modulo
\[
\Delta x_1=0{,}0400\cdot10{,}0=0{,}400\,\mathrm{m}.
\]

Dopo il lancio di \(B\), l'uomo porta ancora \(C\). Il bilancio della quantita
di moto dell'intero sistema e
\[
m_Av_A+m_Bv_B+(M+m_C)u_2=0.
\]
Quindi
\[
u_2=-\frac{1\cdot3+2\cdot2}{70+3}
=-\frac{7}{73}
\simeq-0{,}0959\,\mathrm{m/s}.
\]
Nel secondo intervallo di \(10{,}0\,\mathrm{s}\):
\[
\Delta x_2\simeq0{,}959\,\mathrm{m}.
\]
Al tempo \(t=20{,}0\,\mathrm{s}\) l'uomo ha percorso
\[
x_{20}^{(1)}\simeq0{,}400+0{,}959=1{,}359\,\mathrm{m}.
\]

Dopo il terzo lancio tutte le palle sono separate dall'uomo. La sua velocita
finale vale
\[
u_3=-\frac{m_Av_A+m_Bv_B+m_Cv_C}{M}
=-\frac{3+4+3}{70}
=-0{,}1429\,\mathrm{m/s}.
\]
La distanza residua dal bordo e
\[
5{,}00-1{,}359=3{,}641\,\mathrm{m},
\]
per cui
\[
\Delta t_3^{(1)}=\frac{3{,}641}{0{,}1429}
\simeq25{,}49\,\mathrm{s}.
\]
Il tempo totale e quindi
\[
T_1\simeq45{,}49\,\mathrm{s}.
\]

\textbf{Scenario \(C\)-\(B\)-\(A\).}

Dopo il lancio di \(C\), l'uomo porta ancora \(A\) e \(B\):
\[
w_1=-\frac{m_Cv_C}{M+m_A+m_B}
=-\frac{3}{73}
\simeq-0{,}0411\,\mathrm{m/s}.
\]
Nei primi \(10{,}0\,\mathrm{s}\):
\[
\Delta y_1\simeq0{,}411\,\mathrm{m}.
\]

Dopo il lancio di \(B\), l'uomo porta ancora \(A\):
\[
w_2=-\frac{m_Cv_C+m_Bv_B}{M+m_A}
=-\frac{3+4}{71}
\simeq-0{,}0986\,\mathrm{m/s}.
\]
Nel secondo intervallo percorre
\[
\Delta y_2\simeq0{,}986\,\mathrm{m},
\]
quindi al tempo \(t=20{,}0\,\mathrm{s}\):
\[
x_{20}^{(2)}\simeq1{,}397\,\mathrm{m}.
\]

Dopo l'ultimo lancio la quantita di moto totale delle tre palle e la stessa del
primo scenario, percio
\[
w_3=-\frac{10}{70}
=-0{,}1429\,\mathrm{m/s}.
\]
La distanza residua vale
\[
5{,}00-1{,}397=3{,}603\,\mathrm{m},
\]
e il tempo necessario a percorrerla e
\[
\Delta t_3^{(2)}=\frac{3{,}603}{0{,}1429}
\simeq25{,}22\,\mathrm{s}.
\]
Quindi
\[
T_2\simeq45{,}22\,\mathrm{s}.
\]

Poiche \(T_2<T_1\), l'uomo raggiunge prima il bordo nello scenario
\(C\)-\(B\)-\(A\). La velocita finale, invece, e la stessa nei due casi:
\[
|v_f|=0{,}1429\,\mathrm{m/s}\simeq0{,}143\,\mathrm{m/s}.
\]
L'ordine dei lanci modifica gli spostamenti accumulati nei primi due intervalli,
ma non la quantita di moto totale trasferita alle palle quando tutti i lanci
sono terminati.
\end{soluzione}
